% !TEX root = ../../ctfp-print.tex

\lettrine[lhang=0.17]{是时候}{好好谈谈}集合了. 数学家们对集合论爱恨交加. 它是数学的汇编语言 ---
至少曾经是. 范畴论在某种程度上试图脱离集合论. 比如说, 一个众所周知的事实是: 所有集合的集合并不存在,
但所有集合构成的范畴 $\Set$ 是存在的. 这很好. 另一方面, 我们又假设范畴中任意两个物件之间的态射构成
一个集合. 我们甚至叫它 hom-集. 平心而论, 范畴论确实有一个分支, 那里的态射不构成集合, 而是另一个范畴
里的物件. 这类使用 hom-物件而非 hom-集的范畴叫作\newterm{enriched 富}范畴. 不过在下文中, 我们还是
坚持使用老派的 hom-集范畴.

集合是你所能得到的、最接近于毫无特征的团块的东西 --- 范畴论物件除外. 集合有元素, 但关于这些元素你
说不出太多东西. 如果有一个有限集合, 你可以数它的元素. 对无穷集合的元素, 你可以用基数某种程度地数一数.
比如自然数的集合就比实数的集合小, 尽管两者都是无穷的. 但也许令人意外的是, 有理数的集合与自然数的集合
一样大.

除此之外, 关于集合的全部信息都能编码进集合之间的函数里 --- 尤其是那些可逆的、叫作同构的函数.
就一切实际目的而言, 同构的集合就是相同的. 在我招来基础数学家的怒火之前, 请允许我先解释一下:
相等性与同构之间的区分具有根本的重要性. 事实上, 它是数学最新分支 --- 同伦类型论 (HoTT) --- 的主要
关切之一. 我提到 HoTT, 是因为它是一个从计算中汲取灵感的纯数学理论, 而它的主要倡导者之一
Vladimir Voevodsky, 正是在研究 Coq 定理证明器的过程中获得了那次重大顿悟. 数学与编程的交流是双向的.

关于集合重要的一课是: 比较元素迥异的集合是没有问题的. 例如我们可以说, 某个给定的自然变换的集合与
某个态射的集合同构, 因为集合就是集合而已. 这里的同构只是说, 一边集合中的每个自然变换都对应着另一边
集合中唯一的态射, 反之亦然. 它们可以一一配对. 如果苹果和橙子来自不同的范畴, 你不能拿苹果跟橙子比,
但你可以拿苹果的集合与橙子的集合相比. 常常是把一个范畴论问题转换成集合论问题, 就能给我们所需的洞见,
甚至让我们证明出有价值的定理.

\section{Hom 函子}

每个范畴都自带一族典范的、到 $\Set$ 的映射. 这些映射其实是函子, 所以它们保持范畴的结构.
我们来构造一个这样的映射.

我们在 $\cat{C}$ 中固定一个物件 $a$, 再在 $\cat{C}$ 中挑选另一个物件 $x$.
hom-集 $\cat{C}(a, x)$ 是一个集合, 是 $\Set$ 中的一个物件. 当我们变动
$x$、保持 $a$ 固定时, $\cat{C}(a, x)$ 也会在 $\Set$ 中变动. 于是我们得
到了一个从 $x$ 到 $\Set$ 的映射.

\begin{figure}[H]
  \centering
  \includegraphics[width=0.45\textwidth]{images/hom-set.jpg}
\end{figure}

\noindent
如果想强调我们是把 hom-集看作对它第二个参数的映射, 我们就使用记号
$\cat{C}(a, -)$, 其中的短横线充当参数的占位符.

这个物件间的映射很容易扩展到态射的映射. 取 $\cat{C}$ 中两个任意物件
$x$ 和 $y$ 之间的一个态射 $f$. 在我们刚刚定义的映射下, 物件
$x$ 被映到集合 $\cat{C}(a, x)$, 物件 $y$ 被映到 $\cat{C}(a, y)$.
如果这个映射要成为函子, $f$ 就必须被映到这两个集合之间的一个函数:
$\cat{C}(a, x) \to \cat{C}(a, y)$

我们逐点地定义这个函数, 也就是对每个参数分别定义. 至于参数, 我们应该从
$\cat{C}(a, x)$ 中任取一个元素 --- 叫它 $h$ 吧. 只要首尾相接, 态射就可以组合.
恰好 $h$ 的目标与 $f$ 的源相匹配, 所以它们的组合:
\[f \circ h \Colon a \to y\]
是一个从 $a$ 到 $y$ 的态射. 因此它是 $\cat{C}(a, y)$ 的成员.

\begin{figure}[H]
  \centering
  \includegraphics[width=0.45\textwidth]{images/hom-functor.jpg}
\end{figure}

\noindent
我们刚刚找到了从 $\cat{C}(a, x)$ 到 $\cat{C}(a, y)$ 的函数, 它可以充当
$f$ 的像. 如果没有混淆的危险, 我们就把这个提升后的函数写成: $\cat{C}(a, f)$,
把它对态射 $h$ 的作用写成:
\[\cat{C}(a, f) h = f \circ h\]
由于这个构造在任何范畴里都成立, 它在 Haskell 类型的范畴里也一定成立. 在 Haskell 中,
hom-函子更为人熟知的名字是 \code{Reader} 函子:

\src{snippet01}

\src{snippet02}
现在来考虑: 如果我们不固定 hom-集的源、而是固定它的目标, 会发生什么. 换句话说, 我们在问:
映射 $\cat{C}(-, a)$ 是不是也是个函子. 它是, 但它不是协变的, 而是逆变的. 这是因为同样
这种首尾相接的态射匹配导致的是 $f$ 的后组合 (postcomposition), 而不是像 $\cat{C}(a, -)$
那样的前组合 (precomposition).

我们在 Haskell 里已经见过这个逆变函子. 我们叫它 \code{Op}:

\src{snippet03}

\src{snippet04}
最后, 如果让两个物件都变动, 我们就得到一个前函子 $\cat{C}(-, =)$, 它在第一个参数上是
逆变的, 在第二个参数上是协变的 (为了强调两个参数可以独立变动, 我们用双短横线作为第二个占位符).
我们在谈函子性的时候见过这个前函子:

\src{snippet05}
重要的一课是: 这个观察在任何范畴里都成立: 物件到 hom-集的映射是函子性的. 由于逆变等价于
一个从反范畴出发的映射, 我们可以把这个事实简洁地陈述为:
\[C(-, =) \Colon \cat{C}^{op} \times \cat{C} \to \Set\]

\section{可表函子}

我们已经看到, 对 $\cat{C}$ 中每个 $a$ 的选择, 我们都得到一个从
$\cat{C}$ 到 $\Set$ 的函子. 这种保持结构的、到 $\Set$ 的映射常叫作
\newterm{representation 表示}. 我们把 $\cat{C}$ 的物件和态射表示成
$\Set$ 中的集合和函数.

函子 $\cat{C}(a, -)$ 本身有时也叫作可表的. 更一般地, 任何与 hom-函子自然同构
(对某个 $a$ 的选择) 的函子 $F$ 都叫作\newterm{representable 可表的}.
这样的函子必须取值为 $\Set$, 因为 $\cat{C}(a, -)$ 就是.

我之前说过, 我们常常把同构的集合看作相同的. 更一般地, 我们把范畴中同构的
\emph{物件}看作相同的. 这是因为除了通过态射与其他物件 (以及自身) 的关系之外,
物件没有任何结构.

举个例子, 我们之前谈过幺半群的范畴 $\cat{Mon}$, 它最初就是用集合建模的. 但我们小心地
只挑选那些保持集合之幺半群结构的函数作为态射. 所以如果 $\cat{Mon}$ 中两个物件同构,
也就是它们之间存在一个可逆态射, 它们就有着完全相同的结构. 如果我们偷看它们所基于的集合和函数,
会看到一个幺半群的单位元素被映到另一个的单位元素, 两个元素的积被映到它们的像的积.

同样的推理可以应用于函子. 两个范畴之间的函子构成一个范畴, 其中自然变换扮演着态射的角色.
所以两个函子同构、可以被看作相同的, 当且仅当它们之间存在一个可逆的自然变换.

让我们从这个视角分析可表函子的定义. 要使 $F$ 可表, 我们要求: $\cat{C}$
中存在一个物件 $a$; 一个从 $\cat{C}(a, -)$ 到 $F$ 的自然变换
$\alpha$; 另一个反方向的自然变换 $\beta$; 并且它们的复合是恒等自然变换.

看看 $\alpha$ 在某个物件 $x$ 处的分量. 它是 $\Set$ 中的一个函数:
\[\alpha_x \Colon \cat{C}(a, x) \to F x\]
这个变换的自然性条件告诉我们, 对任意从 $x$ 到 $y$ 的态射 $f$,
下面的图式交换:
\[F f \circ \alpha_x = \alpha_y \circ \cat{C}(a, f)\]
在 Haskell 里, 我们会用多态函数代替自然变换:

\src{snippet06}
带上可选的 \code{forall} 量词. 自然性条件

\src{snippet07}
由于参数化性而自动成立 (它就是我早先提到的那些免费定理之一), 但要明白左边的
\code{fmap} 是由函子 $F$ 定义的, 而右边那个是由 reader 函子定义的. 由于 reader
的 \code{fmap} 不过是函数的前组合, 我们还可以更明确一些. 作用在 $h$ ---
$\cat{C}(a, x)$ 的一个元素 --- 上, 自然性条件简化为:

\src{snippet08}
另一个变换 \code{beta} 朝相反的方向走:

\src{snippet09}
它必须遵守自然性条件, 而且必须是 \code{alpha} 的逆:

\begin{snip}{text}
alpha . beta = id = beta . alpha
\end{snip}
稍后我们会看到, 只要 $F a$ 非空, 从 $\cat{C}(a, -)$ 到任意取值于
$\Set$ 的函子的自然变换总是存在 (米田引理), 只是它未必可逆.

让我用 Haskell 给个例子, 拿列表函子、以 \code{Int} 充当 \code{a}.
这里有一个能干这活的自然变换:

\src{snippet10}
我随意挑了数字 12, 用它创建了一个单例列表. 然后我可以对这个列表 \code{fmap}
函数 \code{h}, 得到一个由 \code{h} 的返回类型构成的列表. (实际上, 这样的变换
有无穷多个, 与整数列表一样多.)

自然性条件等价于 \code{map} (列表版的 \code{fmap}) 的可组合性:

\src{snippet11}
但如果我们想找出逆变换, 就得从任意类型 \code{x} 的列表走到一个返回 \code{x}
的函数:

\src{snippet12}
你也许会想从列表里取出一个 \code{x}, 比如用 \code{head}, 但那样对空列表行不通.
注意, 无论选什么类型 \code{a} (取代 \code{Int}) 在这里都行不通. 所以列表函子
不是可表的.

还记得我们谈 Haskell (自) 函子有点像容器的时候吗? 同理, 我们可以把可表函子看作存储函数调用
的记忆化结果的容器 (在 Haskell 里, hom-集的成员就是函数). 表示者物件 ---
$\cat{C}(a, -)$ 中的类型 $a$ --- 被看作键类型, 凭它可以访问一个函数被列表化的值.
我们叫作 \code{alpha} 的那个变换叫 \code{tabulate}, 它的逆 \code{beta}
叫 \code{index}. 这是一个 (略微简化的) \code{Representable} 类定义:

\src{snippet13}
注意表示者类型 --- 我们的 $a$, 这里叫 \code{Rep f} --- 是 \code{Representable}
定义的一部分. 那个星号只是说 \code{Rep f} 是一个类型 (而不是类型构造子或其他更奇异的
kind).

无穷列表, 即不可能为空的流 (stream), 是可表的.

\src{snippet14}
你可以把它们看作接受 \code{Integer} 参数的函数的记忆化值. (严格来说, 我应该用
非负自然数, 只是不想把代码搞复杂.)

要 \code{tabulate} 这样一个函数, 你就创建一个值的无穷流. 当然, 这只可能实现是因为
Haskell 是惰性的. 这些值按需求值. 你用 \code{index} 访问记忆化的值:

\src{snippet15}
有趣的是, 你可以实现单一的记忆化方案, 覆盖一整族返回类型任意的函数.

逆变函子的可表性也类似地定义, 只不过我们固定的是 $\cat{C}(-, a)$ 的第二个参数.
或者等价地, 我们可以考虑从 $\cat{C}^{op}$ 到 $\Set$ 的函子, 因为
$\cat{C}^{op}(a, -)$ 与 $\cat{C}(-, a)$ 是一回事.

可表性还有一个有趣的花样. 记住, 在笛卡尔闭范畴里, hom-集在内部可以被当作指数物件.
hom-集 $\cat{C}(a, x)$ 等价于 $x^a$, 而对一个可表函子 $F$
我们可以写: $-^a = F$.

就当找乐子, 我们对两边取对数: $a = \mathbf{log}F$

当然, 这只是个纯形式的变换, 但如果你了解对数的某些性质, 它会相当有用. 特别地, 事实表明:
基于积类型的函子可以用余积类型来表示, 而基于余积类型的函子一般不可表 (例子: 列表函子).

最后请注意, 可表函子给我们提供了同一样东西的两种不同实现 --- 一个是函数, 一个是数据结构.
它们的内容完全相同 --- 用相同的键取回相同的值. 这就是我所说的 ``相同'' 的含义. 两个自然同构的
函子, 就内容而言是完全一样的. 另一方面, 这两种表示往往实现方式不同, 性能特征也可能不同.
记忆化是一种性能优化手段, 可能大幅缩短运行时间. 能够为同一个底层计算生成不同的表示, 在实践中
非常有价值. 所以, 令人惊讶的是, 尽管范畴论根本不关心性能, 它却提供了大量机会去探索具有实用
价值的替代实现.

\section{挑战}

\begin{enumerate}
  \tightlist
  \item
        证明 hom-函子把 \emph{C} 中的恒等态射映到 $\Set$ 中相应的恒等函数.
  \item
        证明 \code{Maybe} 不是可表的.
  \item
        \code{Reader} 函子是可表的吗?
  \item
        利用 \code{Stream} 表示, 记忆化一个对自己的参数平方的函数.
  \item
        证明 \code{Stream} 的 \code{tabulate} 和 \code{index} 确确实实互逆.
        (提示: 用归纳法.)
  \item
        函子:

        \begin{snip}{haskell}
Pair a = Pair a a
\end{snip}
        是可表的. 你能猜出表示它的类型吗? 实现 \code{tabulate} 和 \code{index}.
\end{enumerate}

\section{参考资料}

\begin{enumerate}
  \tightlist
  \item
        The Catsters 关于
        \urlref{https://www.youtube.com/watch?v=4QgjKUzyrhM}{可表函子}的视频.
\end{enumerate}
